% Original paper introduction — commented out because the pedagogical
% introduction in Sections 8.0–8.3 of this chapter covers all the same
% material in greater depth.
%
% \subsection{Introduction}\label{sec:Intro}
%
% {High-energy cosmic messengers offer a glimpse into some of the most powerful phenomena in the Universe. Among these messengers, $\gamma$-rays stand out as a unique tool for exploring the physical processes driving astrophysical sources such as active galactic nuclei (AGN), supernovae, star-forming galaxies, and pulsars. At the same time, $\gamma$-rays are among the most promising observational signatures expected from dark matter (DM), should it consist of a new particle capable of self-annihilation or decay.}
%
% The Cherenkov Telescope Array Observatory (CTAO) \cite{CTAConsortium:2018tzg, mazin2019cherenkov}, designed to detect cosmic $\gamma$-rays with energies above 20 GeV, will offer a unique opportunity to probe both the high-energy emission of astrophysical sources and potential signals from DM particles with masses in the TeV range. CTAO will improve upon the performance of current imaging atmospheric Cherenkov telescopes, thanks to its better angular resolution, larger field of view, and broader energy coverage. It will also be the first ground-based Cherenkov telescope to carry out an intensive program of dedicated sky surveys {\cite{CTAConsortium:2018tzg, mazin2019cherenkov}}, covering large regions (i.e.\ more than a quarter) of the sky. These capabilities make CTAO particularly well-suited to studying the cosmological $\gamma$-ray emission potentially produced by DM annihilation or decay. However, this emission is irreducibly intermingled with the non-thermal emissions from astrophysical sources. The diffuse unresolved $\gamma$-ray background \cite{Fornasa_2015} (UGRB) consists of the cumulative emission from $\gamma$-ray astrophysical sources which are too faint to be individually resolved, with a possible additional contribution from DM \cite{Ajello:2015mfa}.
%
% A potentially powerful method to separate the astrophysical emission from a DM contribution is the analysis of anisotropies in the $\gamma$-ray sky.
% Various techniques are possible for this purpose.
% The present work focuses in particular on the auto-correlation of the gamma-ray signal and
% cross-correlations between the $\gamma$-ray emission and a gravitational tracer of the DM distribution in the Universe.
% Auto-correlation analysis are especially well suited for studying the properties of rare and bright astrophysical populations, like blazars, as opposed to faint and numerous populations, such as star-forming galaxies \cite{Ando:2006mt,Ando:2013ff,Fermi-LAT:2018udj,fornasa2016angular,Korsmeier:2022cwp}.
% Anisotropy auto-correlation studies specifically in the TeV band for Cherenkov telescopes have been performed in Refs. \cite{Hutten:2018wop,Hutten:2016jko,Ripken:2012db}.
% The cross-correlation technique, instead, offers the advantage of incorporating direct gravitational tracers of the matter distribution in the Universe, and can leverage differences in angular scale, energy spectrum, and redshift evolution between the cross-correlated fields \cite{camera2013novel,camera2015tomographic,fornengo2014particle}. This approach provides a multi-dimensional handle to disentangle signals with distinct features in their spatial distribution, energy spectra, and redshift dependence. Extragalactic astrophysical sources, for instance, typically appear as point-like emitters at $\gamma$-ray energies, whereas DM is expected to produce a diffuse signal that traces the large-scale structure of the Universe. Different classes of astrophysical sources exhibit diverse spectral energy distributions, which are generally distinct from those predicted for DM-induced emissions. Moreover, their redshift distributions differ: DM $\gamma$-ray emission is expected to peak at low redshift, unlike unresolved astrophysical populations {\cite{camera2013novel,camera2015tomographic,fornengo2014particle}}. As a result, DM signals are expected to correlate more strongly with low-redshift gravitational tracers.
% Beyond its application in the search for DM, the cross-correlation technique also holds broader potential for studying unresolved $\gamma$-ray source populations, offering insights into their redshift distribution and clustering properties.
%
% Cross-correlation analyses have been performed between \textit{Fermi}-LAT $\gamma$-ray sky maps and various tracers of the large-scale structure. These include weak gravitational lensing  \cite{shirasaki2014cross,shirasaki2016cosmological,troster2017cross,DES:2019ucp}, the clustering of galaxies \cite{xia2015tomography,regis2015particle,cuoco2015dark,shirasaki2015cross,cuoco2017tomographic,Ammazzalorso:2018evf,paopiamsap2024constraints}, galaxy clusters \cite{branchini2017cross,shirasaki2018correlation,hashimoto2019measurement,colavincenzo2020searching,tan2020bounds}, and the lensing effect of the cosmic microwave background \cite{fornengo2015evidence}, which traces
% the large-scale distribution of matter across cosmological distances. Additional studies are presented in Refs. \cite{ando2014mapping,fornengo2014particle,fornasa2016angular,feng2017planck,Tan:2020fbc}).
%
% Positive cross-correlation signals have been found with galaxies at the level of $3.5\sigma$ \cite{xia2015tomography,regis2015particle,cuoco2015dark} up to $8-10\sigma$ \cite{paopiamsap2024constraints}, galaxy clusters at $4.7\sigma$ \cite{branchini2017cross}, lensing of the cosmic microwave background at $3.2\sigma$ \cite{fornengo2015evidence} and weak gravitational lensing at $5.3\sigma$ \cite{DES:2019ucp} and more recently at signal-to-noise ratio of 8.9  \cite{DES:2025ulp}.
%
% In this work, we explore the potential of the cross-correlation and auto-correlation techniques to detect a signal originating from sub-threshold extragalactic source populations or from annihilation or decay of TeV-scale DM particles by using the CTAO extragalactic survey capabilities.

% Bridge paragraph connecting the chapter introduction to the paper body
In this paper, we present a sensitivity forecast for detecting dark matter signals through cross-correlations between the CTAO extragalactic survey and the 2MASS galaxy catalog. Section~\ref{sec:CTA} describes the CTAO instrument configuration and survey strategy. Section~\ref{sec:formalism} develops the cross-correlation signal model, including the window functions and angular power spectra for both astrophysical sources and dark matter. Section~\ref{sec:sens} presents the sensitivity forecast results for dark matter annihilation and decay. We conclude in Section~\ref{sec:xcorr:conclusions}.

%%%%%%%%%%%%%%%%
